\sectioncentered*{Заключение}
\addcontentsline{toc}{section}{Заключение}

Цель курсового проекта достигнута: спроектирована и реализована автоматизированная информационная система отслеживания цен на товары в виде \textit{Telegram}-бота. Система каждые 3 минуты опрашивает площадку \textit{Kufar.by}, отбирает объявления дешевле заданной пользователем цены и уведомляет о новых предложениях и о снижении цены без повторов.

Анализ предметной области показывает, что ручной мониторинг не успевает за круглосуточным появлением объявлений, а существующие средства решают лишь часть задачи: сохранённые поиски площадки не отслеживают цену во времени, зарубежные трекеры рассчитаны на новые товары с артикулом, открытые боты ограничиваются текстовым поиском.

Ядро работы~-- база данных под управлением \textit{PostgreSQL}. \textit{ER}-модель в нотации Питера Чена из семи сущностей переходит в логическую и физическую схемы. Схема отвечает третьей нормальной форме, а отступления сделаны осознанно: \textit{JSON}-поля хранят фильтры категорий, журнал уведомлений~-- цену на момент отправки для проверки повторов. Индексы подобраны под запросы цикла опроса и доставки, прежде всего под поиск объявления по внешнему идентификатору~-- более 500 тысяч раз в сутки. Целостность поддерживают внешние ключи, ограничения уникальности и перечислимые типы.

Программная часть разделена на серверный сервис на \textit{FastAPI}, который выполняет SQL-запросы и периодические задачи, и бота на \textit{aiogram} для диалога с пользователем. Отбор по бренду, размеру рамы или числу комнат выполняет сама площадка. Система реализует весь функционал из постановки задачи: импорт фильтров из ссылки, мгновенный поиск, обмен списком товаров через \textit{CSV}, отчёты в \textit{Excel}, два языка интерфейса и копии уведомлений на почту.

Надёжность опирается на свойства СУБД. Найденное совпадение сначала фиксируется в базе данных, а отправляет его отдельная задача с повтором до пяти попыток; отметка о доставке и запись в журнал сохраняются одной транзакцией. Работу системы подтверждают более 280 автоматических тестов.

Система допускает развитие: модель данных предусматривает несколько площадок, поэтому к \textit{Kufar.by} добавляются \textit{Onliner}, \textit{av.by} и другие площадки. Система решает реальную задачу покупателей и перекупщиков: внешне это бот с кнопками, по существу~-- система обработки потока данных, качество которой определяет проектирование базы данных.
